
\documentclass[a4paper,12pt]{article}

\usepackage[slovene]{babel}
\usepackage[utf8]{inputenc}
\usepackage[T1]{fontenc}
\usepackage{lmodern}
\usepackage{amsmath,amssymb,amsfonts}
\usepackage{mathtools}

\setlength{\parindent}{0mm}

\begin{document}

\section*{Logika in množice}

\begin{enumerate}
    \item
    Poišči preneksno obliko formule
    $\exists x:P(x)\land\forall x:Q(x)\Rightarrow\forall x:R(x)$.

    \item
    Definiramo množici $A=[2,5]$ in
    $B=\{0,1,2,3,4,\ldots\}$.
    V ravnino nariši:
    \begin{enumerate}
        \item $A\cap B\times\emptyset$
        \item $(A\cup B)^c\times\mathbb{R}$
    \end{enumerate}

    \item
    Dokaži:
    \begin{itemize}
        \item $(A\Rightarrow B)\sim(\neg B\Rightarrow\neg A)$
        \item $\neg(A\lor B)\sim\neg A\land\neg B$
    \end{itemize}
\end{enumerate}

\section*{Analiza}

\begin{enumerate}
    \item
    Pokaži, da je funkcija
    $x\mapsto\sqrt{x}$
    enakomerno zvezna na $[0,\infty)$.

    \item
    Katero krivuljo določa naslednji parametrični zapis?
    \[
    x(t)=a\cos t,\qquad
    y(t)=b\sin t,\qquad
    t\in[0,2\pi]
    \]

    \item
    Pokaži, da ima
    $f(x)=3x+\sin(2x)$
    inverzno funkcijo in izračunaj
    $(f^{-1})'(3\pi)$.

    \item
    Izračunaj integral
    \[
    \int\frac{2+\sqrt{x+1}}
    {(x+1)^2-\sqrt{x+1}}\,dx
    \]

    \item
    Naj bo $g$ zvezna funkcija.
    Ali posplošeni integral
    $\int_0^1\frac{g(x)}{x^2}\,dx$
    konvergira ali divergira?
    Utemelji.

    \item
    S Taylorjevo formulo pokaži, da v okolici točke
    $0$ velja
    \[
    \sin x=x-\frac{x^3}{3!}+O(x^5).
    \]
\end{enumerate}

\section*{Kompleksna števila}

\begin{enumerate}
    \item
    Naj bo $z$ kompleksno število,
    $z\ne1$ in $|z|=1$.
    Dokaži, da je število
    $i\frac{z+1}{z-1}$ realno.

    \item
    Poenostavi izraz:
    \[
    \frac{
        \dfrac{3+i}{2-2i}+\dfrac{7i}{1-i}
    }{
        1+\dfrac{i-1}{4}-\dfrac{5}{2-3i}
    }
    \]
\end{enumerate}

\section*{Zaporedja, vrste in limite}

\begin{enumerate}
    \item
    Z indukcijo dokaži, da za vsako naravno število
    $n$ velja
    \[
    \sum_{i=1}^{n}i^2=
    \frac{n(n+1)(2n+1)}{6}.
    \]

    \item
    Naj bo $\sum_{n=1}^{\infty}a_n$
    absolutno konvergentna vrsta in
    $a_n\ne-1$.
    Dokaži, da je tudi vrsta
    $\sum_{n=1}^{\infty}\frac{a_n}{1+a_n}$
    absolutno konvergentna.

    \item
    Izračunaj limito
    \[
    \lim_{x\to\infty}
    \left(\sin\sqrt{x+1}-\sin\sqrt{x}\right).
    \]

    \item
    Za dani zaporedji preveri, ali sta konvergentni.
    \[
    a_n=
    \underbrace{
    \sqrt{2+\sqrt{2+\cdots+\sqrt{2}}}
    }_{\text{\(n\) korenov}}
    \qquad
    b_n=
    \underbrace{
    \sin(\sin(\ldots(\sin1)\ldots))
    }_{\text{\(n\) sinusov}}
    \]
\end{enumerate}

\section*{Algebra}

\begin{enumerate}
    \item
    Vektorja
    $\vec{c}=\vec{a}+2\vec{b}$
    in
    $\vec{d}=\vec{a}-\vec{b}$
    sta pravokotna in imata dolžino 1.
    Določi kot med vektorjema
    $\vec{a}$ in $\vec{b}$.

    \item
    Izračunaj
    \[
    \begin{pmatrix}
    1&2&3&4&5&6\\
    4&5&2&6&3&1
    \end{pmatrix}^{-2000}
    \]
\end{enumerate}

\section*{Velika determinanta}

Izračunaj naslednjo determinanto
$2n\times2n$, ki ima na neoznačenih mestih ničle.

\[
\left|
\begin{array}{ccccccccc}
1&&&1&&&&&\\
&2&&2&&&&&\\
&&\ddots&\vdots&&&&&\\
&&n-1&n-1&&&&&\\
1&2&\cdots&n-1&n&n+1&n+2&\cdots&2n\\
&&&n+1&n+1&&&&\\
&&&n+2&&n+2&&&\\
&&&\vdots&&&\ddots&&\\
&&&2n&&&&&2n
\end{array}
\right|
\]

\section*{Grupe}

Naj bo
\begin{align*}
G={}&\{z\in\mathbb{C};
z=2^k(\cos(m\pi\sqrt{2})+
i\sin(m\pi\sqrt{2})),\
k,m\in\mathbb{Z}\}\\
H={}&\{(x,y)\in\mathbb{R}^2;\
x,y\in\mathbb{Z}\}
\end{align*}

\begin{enumerate}
    \item
    Pokaži, da je $G$ podgrupa v grupi
    $(\mathbb{C}\setminus\{0\},\cdot)$
    neničelnih kompleksnih števil za običajno množenje.

    \item
    Pokaži, da je $H$ podgrupa v aditivni grupi
    $(\mathbb{R}^2,+)$ ravninskih vektorjev
    za običajno seštevanje po komponentah.

    \item
    Pokaži, da je preslikava
    $f:H\to G$, podana s pravilom
    \[
    (x,y)\mapsto
    2^x\bigl(\cos(y\pi\sqrt{2})+
    i\sin(y\pi\sqrt{2})\bigr)
    \]
    izomorfizem grup $G$ in $H$.
\end{enumerate}

\end{document}
